% GENERATED FILE. Edit canonical lessons, then run npm run generate:books.
% canonical-source-hash: fnv1a64:9056148822ea9e9d
% canonical-lessons: HI-C454-cakit, HI-C454-ascarya, HI-C454-acambha, HI-C454-ajib, HI-C454-kamal

\chapter{A Surprise}
\label{ch:hi-a-surprise}
\hlchaptermodality{454}

\begin{chapteropening}
\textbf{By the end of this chapter:} \emph{I can say that something surprised me.}

Complete the last lesson of chapter 454: I can say that something surprised me.
\end{chapteropening}

\section[\texorpdfstring{\hi{चकित} (cakit)}{चकित (cakit)}]{\hi{चकित} (cakit) — taken aback; astonished}

\label{lesson:HI-C454-cakit}



\begin{quote}
Before the new one: say the Hindi for elation, then the Hindi for pleased.
\end{quote}

\subsection*{You'll want to know: \hi{चकित}}
\textbf{\hi{चकित}} — \emph{cakit} — \textquotedblleft{}taken aback; astonished\textquotedblright{}.

\begin{quote}
\hi{मीरा} \hi{का} \hi{नया} \hi{घर} \hi{देखकर} \hi{सब} \hi{चकित} \hi{थे।}
\end{quote}

\begin{quote}
\emph{Mīrā kā nayā ghar dekhkar sab cakit the.}
\end{quote}

\begin{quote}
Everyone was astonished when they saw Mira's new house.
\end{quote}

\begin{cousinweb}
\textbf{\hi{चकित}} (\emph{cakit}) is a tatsama, borrowed from Sanskrit \emph{cakita}, \textquotedblleft{}trembling, timid, startled\textquotedblright{} (Monier-Williams). Hindi has turned the start of fright into the start of surprise.
\end{cousinweb}

\begin{grammarlens}[title={What you've built}]
One more word for describing what you lived and felt.
\end{grammarlens}

\subsection*{Your turn}
\begin{itemize}
\raggedright
  \item \emph{Say it:} \emph{cakit}
  \item \emph{Say it:} \emph{cakit}, once more
  \item \emph{Say it:} \emph{Mīrā kā nayā ghar dekhkar sab cakit the.} — Everyone was astonished when they saw Mira's new house.
  \item \emph{Recall:} say the Hindi for elation, then the Hindi for pleased, then say \emph{cakit} again
\end{itemize}

\begin{tcolorbox}[breakable,colback=teal!4,colframe=teal!35!black,title={Before you move on}]
What does \hi{चकित} mean? (\textquotedblleft{}Taken aback; astonished\textquotedblright{}.) Say it once more.
\end{tcolorbox}
\section[\texorpdfstring{\hi{आश्चर्य} (āścarya)}{आश्चर्य (āścarya)}]{\hi{आश्चर्य} (āścarya) — surprise; wonder}

\label{lesson:HI-C454-ascarya}



\begin{quote}
Before the new one: say the Hindi for pleased, then the Hindi for taken aback.
\end{quote}

\subsection*{You'll want to know: \hi{आश्चर्य}}
\textbf{\hi{आश्चर्य}} — \emph{āścarya} — \textquotedblleft{}surprise; wonder\textquotedblright{}.

\begin{quote}
\hi{मुझे} \hi{अरुण} \hi{की} \hi{बात} \hi{पर} \hi{बहुत} \hi{आश्चर्य} \hi{हुआ।}
\end{quote}

\begin{quote}
\emph{mujhe Aruṇ kī bāt par bahut āścarya huā.}
\end{quote}

\begin{quote}
I was very surprised at what Arun said.
\end{quote}

Masculine.

\begin{cousinweb}
\textbf{\hi{आश्चर्य}} (\emph{āścarya}) is a tatsama, borrowed from Sanskrit \emph{āścarya}, \textquotedblleft{}appearing rarely, wonderful\textquotedblright{}, and as a noun \textquotedblleft{}a wonder\textquotedblright{} (Monier-Williams).
\end{cousinweb}

\begin{grammarlens}[title={What you've built}]
One more word for describing what you lived and felt.
\end{grammarlens}

\subsection*{Your turn}
\begin{itemize}
\raggedright
  \item \emph{Say it:} \emph{āścarya}
  \item \emph{Say it:} \emph{āścarya}, once more
  \item \emph{Say it:} \emph{mujhe Aruṇ kī bāt par bahut āścarya huā.} — I was very surprised at what Arun said.
  \item \emph{Recall:} say the Hindi for pleased, then the Hindi for taken aback, then say \emph{āścarya} again
\end{itemize}

\begin{tcolorbox}[breakable,colback=teal!4,colframe=teal!35!black,title={Before you move on}]
What does \hi{आश्चर्य} mean? (\textquotedblleft{}Surprise; wonder\textquotedblright{}.) Say it once more.
\end{tcolorbox}
\section[\texorpdfstring{\hi{अचंभा} (acambhā)}{अचंभा (acambhā)}]{\hi{अचंभा} (acambhā) — astonishment; a marvel}

\label{lesson:HI-C454-acambha}



\begin{quote}
Before the new one: say the Hindi for taken aback, then the Hindi for surprise.
\end{quote}

\subsection*{You'll want to know: \hi{अचंभा}}
\textbf{\hi{अचंभा}} — \emph{acambhā} — \textquotedblleft{}astonishment; a marvel\textquotedblright{}.

\begin{quote}
\hi{यह} \hi{अचंभे} \hi{की} \hi{बात} \hi{है} \hi{कि} \hi{दादा} \hi{अस्सी} \hi{साल} \hi{के} \hi{हैं} \hi{और} \hi{रोज़} \hi{दौड़ते} \hi{हैं।}
\end{quote}

\begin{quote}
\emph{yah acambhe kī bāt hai ki dādā assī sāl ke haiṁ aur roz dauṛte haiṁ.}
\end{quote}

\begin{quote}
It's a marvel that my grandfather is eighty and still runs every day.
\end{quote}

Masculine.

\begin{cousinweb}
\textbf{\hi{अचंभा}} (\emph{acambhā}) is a tadbhava: CDIAL lists Hindi \emph{acambhā}, \textquotedblleft{}wonderful; wonder\textquotedblright{}, under \emph{*ācchambha}, a form it reconstructs (CDIAL 1079).
\end{cousinweb}

\begin{grammarlens}[title={What you've built}]
One more word for describing what you lived and felt.
\end{grammarlens}

\subsection*{Your turn}
\begin{itemize}
\raggedright
  \item \emph{Say it:} \emph{acambhā}
  \item \emph{Say it:} \emph{acambhā}, once more
  \item \emph{Say it:} \emph{yah acambhe kī bāt hai ki dādā assī sāl ke haiṁ aur roz dauṛte haiṁ.} — It's a marvel that my grandfather is eighty and still runs every day.
  \item \emph{Recall:} say the Hindi for taken aback, then the Hindi for surprise, then say \emph{acambhā} again
\end{itemize}

\begin{tcolorbox}[breakable,colback=teal!4,colframe=teal!35!black,title={Before you move on}]
What does \hi{अचंभा} mean? (\textquotedblleft{}Astonishment; a marvel\textquotedblright{}.) Say it once more.
\end{tcolorbox}
\section[\texorpdfstring{\hi{अजीब} (ajīb)}{अजीब (ajīb)}]{\hi{अजीब} (ajīb) — strange; odd}

\label{lesson:HI-C454-ajib}



\begin{quote}
Before the new one: say the Hindi for surprise, then the Hindi for astonishment.
\end{quote}

\subsection*{You'll want to know: \hi{अजीब}}
\textbf{\hi{अजीब}} — \emph{ajīb} — \textquotedblleft{}strange; odd\textquotedblright{}.

\begin{quote}
\hi{रात} \hi{को} \hi{खिड़की} \hi{से} \hi{एक} \hi{अजीब} \hi{आवाज़} \hi{आई।}
\end{quote}

\begin{quote}
\emph{rāt ko khiṛkī se ek ajīb āvāz āī.}
\end{quote}

\begin{quote}
A strange sound came through the window at night.
\end{quote}

\begin{cousinweb}
\textbf{\hi{अजीब}} (\emph{ajīb}) is usually traced to Arabic \emph{ʿajīb}, \textquotedblleft{}wonderful, strange\textquotedblright{}, which reached Hindi through Persian.
\end{cousinweb}

\begin{grammarlens}[title={What you've built}]
One more word for describing what you lived and felt.
\end{grammarlens}

\subsection*{Your turn}
\begin{itemize}
\raggedright
  \item \emph{Say it:} \emph{ajīb}
  \item \emph{Say it:} \emph{ajīb}, once more
  \item \emph{Say it:} \emph{rāt ko khiṛkī se ek ajīb āvāz āī.} — A strange sound came through the window at night.
  \item \emph{Recall:} say the Hindi for surprise, then the Hindi for astonishment, then say \emph{ajīb} again
\end{itemize}

\begin{tcolorbox}[breakable,colback=teal!4,colframe=teal!35!black,title={Before you move on}]
What does \hi{अजीब} mean? (\textquotedblleft{}Strange; odd\textquotedblright{}.) Say it once more.
\end{tcolorbox}
\section[\texorpdfstring{\hi{कमाल} (kamāl)}{कमाल (kamāl)}]{\hi{कमाल} (kamāl) — a marvel; something amazing}

\label{lesson:HI-C454-kamal}



\begin{quote}
Before the new one: say the Hindi for astonishment, then the Hindi for strange.
\end{quote}

\subsection*{You'll want to know: \hi{कमाल}}
\textbf{\hi{कमाल}} — \emph{kamāl} — \textquotedblleft{}a marvel; something amazing\textquotedblright{}.

\begin{quote}
\hi{मीरा} \hi{ने} \hi{एक} \hi{दिन} \hi{में} \hi{पूरी} \hi{किताब} \hi{पढ़} \hi{ली}, \hi{कमाल} \hi{है}!
\end{quote}

\begin{quote}
\emph{Mīrā ne ek din meṁ pūrī kitāb paṛh lī, kamāl hai!}
\end{quote}

\begin{quote}
Mira read the whole book in a day — amazing!
\end{quote}

Masculine.

\begin{cousinweb}
\textbf{\hi{कमाल}} (\emph{kamāl}) is usually traced to Arabic \emph{kamāl}, \textquotedblleft{}perfection\textquotedblright{}, which reached Hindi through Persian. \emph{kamāl hai!} is \textquotedblleft{}that's amazing!\textquotedblright{}.
\end{cousinweb}

\begin{grammarlens}[title={What you've built}]
One more word for describing what you lived and felt.
\end{grammarlens}

\subsection*{Your turn}
\begin{itemize}
\raggedright
  \item \emph{Say it:} \emph{kamāl}
  \item \emph{Say it:} \emph{kamāl}, once more
  \item \emph{Say it:} \emph{Mīrā ne ek din meṁ pūrī kitāb paṛh lī, kamāl hai!} — Mira read the whole book in a day — amazing!
  \item \emph{Recall:} say the Hindi for astonishment, then the Hindi for strange, then say \emph{kamāl} again
\end{itemize}

\begin{tcolorbox}[breakable,colback=teal!4,colframe=teal!35!black,title={Before you move on}]
What does \hi{कमाल} mean? (\textquotedblleft{}A marvel; something amazing\textquotedblright{}.) Say it once more.
\end{tcolorbox}
